\newpage

\section{ইসমে মুসতাক}

% ─── Title ──────────────────────────────────────────────────
\begin{center}
  {\arabicfont \color{titlegreen} \textbf{\textarabic{الاسم المشتق}}}
\end{center}

% ─── Definition ─────────────────────────────────────────────
\noindent যে {\arabicfont \textarabic{اسم}} কোনো {\arabicfont \textarabic{مصدر}}
থেকে বের হয়ে বিভিন্ন ওজনে আসে, তাকে {\arabicfont \textarabic{اسم مشتق}} বলে।

\begin{center}
  যেমন: {\arabicfont \textarabic{الْعِلْمُ}} (জ্ঞান) থেকে
  $\rightarrow$
  {\arabicfont \ismfail{ع}{ل}{م}} জ্ঞানী,
  \enspace
  {\arabicfont \ismmaful{ع}{ل}{م}} জ্ঞাত,
  \enspace
  {\arabicfont \ismtafdil{ع}{ل}{م}} অধিক জ্ঞানী
\end{center}

\noindent {\arabicfont \textarabic{اسم مشتق}} ৫ প্রকার:

\vspace{0.2cm}

% ─── Five types table ───────────────────────────────────────
% Grouped so the vertically-centred (m) X columns don't leak into
% other chapters' tables.
\begingroup
\renewcommand{\tabularxcolumn}[1]{m{#1}}
\setlength{\arrayrulewidth}{0.5pt}
\setlength{\tabcolsep}{3pt}
\renewcommand{\arraystretch}{1.6}

\noindent\begin{tabularx}{\textwidth}{|C|C|C|C|}
  \hline
  \rowcolor{headergreen}
  {\color{white} \textbf{প্রকার}}    &
  {\color{white} \textbf{যা বোঝায়}} &
  {\color{white} \textbf{ওজন}}      &
  {\color{white} \textbf{উদাহরণ}}   \\
  \hline

  \rowcolor{lightgreen}
  {\arabicfont \textarabic{اسم فاعل}} &
  কাজের কর্তা &
  {\arabicfont \ismfail{ف}{ع}{ل}} &
  {\arabicfont \ismfail{ن}{ص}{ر}}\newline{\small সাহায্যকারী} \\
  \hline

  {\arabicfont \textarabic{اسم مفعول}} &
  যার উপর কাজ হয় &
  {\arabicfont \ismmaful{ف}{ع}{ل}} &
  {\arabicfont \ismmaful{ن}{ص}{ر}}\newline{\small সাহায্যপ্রাপ্ত} \\
  \hline

  \rowcolor{lightgreen}
  {\arabicfont \textarabic{اسم ظرف}} &
  কাজের সময় বা স্থান &
  {\arabicfont \ismzarf{ف}{ع}{ل}}\newline{\arabicfont \ismzarfi{ف}{ع}{ل}} &
  {\arabicfont \ismzarfi{س}{ج}{د}}\newline{\small সিজদার স্থান} \\
  \hline

  % اسم آلة: three sub-rows (صغرى / وسطى / كبرى); name sits in the middle one.
  & &
  {\arabicfont \ismaalas{ف}{ع}{ل}}\newline
  {\footnotesize ছোট ({\arabicfont \textarabic{صغرى}})} &
  {\arabicfont \ismaalas{ن}{ج}{ل}}\newline{\small কাস্তে} \\
  \cline{3-4}

  {\arabicfont \textarabic{اسم آلة}} &
  কাজের যন্ত্র &
  {\arabicfont \ismaalaw{ف}{ع}{ل}}\newline
  {\footnotesize মাঝারি ({\arabicfont \textarabic{وسطى}})} &
  {\arabicfont \ismaalaw{ك}{ن}{س}}\newline{\small ঝাড়ু} \\
  \cline{3-4}

  & &
  {\arabicfont \ismaalak{ف}{ع}{ل}}\newline
  {\footnotesize বড় ({\arabicfont \textarabic{كبرى}})} &
  {\arabicfont \ismaalak{ف}{ت}{ح}}\newline{\small চাবি} \\
  \hline

  \rowcolor{lightgreen}
  {\arabicfont \textarabic{اسم تفضيل}} &
  তুলনায় বেশি বোঝায় &
  {\arabicfont \ismtafdil{ف}{ع}{ل}} &
  {\arabicfont \ismtafdil{ك}{ب}{ر}}\newline{\small অধিক বড়} \\
  \hline
\end{tabularx}
\endgroup

\vspace{0.4cm}

\begin{center}
  {\arabicfont \color{gray} \small
    \textarabic{الحروف الأصلية بالأسود \textbullet\ الزوائد بالرمادي}
  }
\end{center}

% ════════════════════════════════════════════════════════════
% গঠন প্রণালী — how each form is built from the مضارع
% ════════════════════════════════════════════════════════════
\newpage

\begin{center}
  {\Large \color{titlegreen} গঠন প্রণালী}
\end{center}

\noindent সব {\arabicfont \textarabic{اسم مشتق}} মুজারে থেকে বানানো হয়।
প্রথম ধাপ সবার এক: মুজারের {\arabicfont \textarabic{\sfx{علامة}}} ফেলে দাও, তারপর:

\vspace{0.2cm}

% ─── Construction table ─────────────────────────────────────
% Grouped so the vertically-centred (m) X columns don't leak into
% other chapters' tables.
\begingroup
\renewcommand{\tabularxcolumn}[1]{m{#1}}
\setlength{\arrayrulewidth}{0.5pt}
\setlength{\tabcolsep}{4pt}
\renewcommand{\arraystretch}{1.9}

\noindent\begin{tabularx}{\textwidth}{%
    |>{\hsize=0.62\hsize}C%
    |>{\hsize=0.98\hsize}C%
    |>{\hsize=1.4\hsize\raggedright\arraybackslash}X|}
  \hline
  \rowcolor{headergreen}
  {\color{white} \textbf{প্রকার}} &
  {\color{white} \textbf{গঠন}}    &
  \centering\arraybackslash {\color{white} \textbf{নিয়ম}} \\
  \hline

  \rowcolor{lightgreen}
  {\arabicfont \textarabic{اسم فاعل}} &
  {\arabicfont \mudsrc{يَ}{نْصُرُ}} $\rightarrow$ {\arabicfont \gnaasir} &
  ফা-তে \sfx{যবর}, ফা ও আইনের মাঝে \aug{আলিফ},
  আইনে \sfx{যের}, লামে \sfx{তানবীন} \\
  \hline

  {\arabicfont \textarabic{اسم مفعول}} &
  {\arabicfont \mudsrc{يُ}{نْصَرُ}} $\rightarrow$ {\arabicfont \gmansuur} &
  (মাজহূল থেকে) শুরুতে \sfx{যবর}সহ \aug{মীম}, আইনে \sfx{পেশ},
  আইন ও লামের মাঝে \aug{ওয়াও}, লামে \sfx{তানবীন} \\
  \hline

  \rowcolor{lightgreen}
  {\arabicfont \textarabic{اسم ظرف}} &
  {\arabicfont \mudsrc{يَ}{نْصُرُ}} $\rightarrow$ {\arabicfont \gmansar}\newline
  {\arabicfont \mudsrc{يَ}{ضْرِبُ}} $\rightarrow$ {\arabicfont \gmadrib} &
  শুরুতে \sfx{যবর}সহ \aug{মীম}, আইনে \sfx{যবর}
  (মুজারেতে যের থাকলে যের-ই), লামে \sfx{তানবীন} \\
  \hline

  % اسم آلة: three sub-rows. صغرى comes from the مضارع; وسطى and كبرى
  % are built from the صغرى (plain, uncoloured source). Name in the middle.
  &
  {\arabicfont \mudsrc{يَ}{فْتَحُ}} $\rightarrow$ {\arabicfont \gmiftah} &
  {\arabicfont \footnotesize \textarabic{صغرى}} (ছোট): শুরুতে \sfx{যের}সহ
  \aug{মীম}, আইনে যবর না থাকলে \sfx{যবর}, লামে \sfx{তানবীন} \\
  \cline{2-3}

  {\arabicfont \textarabic{اسم آلة}} &
  {\arabicfont \ismaalas{ف}{ت}{ح}} $\rightarrow$ {\arabicfont \gmiftaha} &
  {\arabicfont \footnotesize \textarabic{وسطى}} (মাঝারি):
  {\arabicfont \footnotesize \textarabic{صغرى}}-র শেষে \aug{গোল~তা},
  লামে \sfx{যবর} \\
  \cline{2-3}

  &
  {\arabicfont \ismaalas{ف}{ت}{ح}} $\rightarrow$ {\arabicfont \gmiftaah} &
  {\arabicfont \footnotesize \textarabic{كبرى}} (বড়):
  {\arabicfont \footnotesize \textarabic{صغرى}}-র আইনের পরে \aug{আলিফ} \\
  \hline

  \rowcolor{lightgreen}
  {\arabicfont \textarabic{اسم تفضيل}} &
  {\arabicfont \mudsrc{يَ}{كْبُرُ}} $\rightarrow$ {\arabicfont \gakbar}\newline
  {\arabicfont \gkubraa} &
  শুরুতে \sfx{যবর}সহ \aug{হামযা}, আইনে \sfx{যবর} (পুং)।
  স্ত্রী: ফা-তে \sfx{পেশ}, আইনে \sfx{সাকিন},
  লামে \sfx{যবর} ও \aug{আলিফ মাকসূরা} \\
  \hline
\end{tabularx}
\endgroup

\vspace{0.4cm}

% ─── مزيد فيه shortcut ───────────────────────────────────────
\noindent{\arabicfont \color{titlegreen} \textbf{\textarabic{مزيد فيه}}} বাবে
মুজারের {\arabicfont \textarabic{\sfx{علامة}}} ফেলে শুরুতে \sfx{পেশ}সহ \aug{মীম} দাও।\\
শেষ হরফের আগে যের $\rightarrow$ \mbox{\arabicfont \textarabic{اسم فاعل}},
যবর $\rightarrow$ \mbox{\arabicfont \textarabic{اسم مفعول}}।

\begin{center}
  {\arabicfont \mazid{يَ}{سْتَ}{عْمِلُ}}
  $\rightarrow$
  {\arabicfont \gmustamil} ব্যবহারকারী,
  \enspace
  {\arabicfont \gmustamal} ব্যবহৃত
\end{center}

\vspace{0.2cm}

\begin{center}
  \small
  কালো = মূল অক্ষর
  \enspace$\bullet$\enspace
  \aug{ধূসর = বাড়তি অক্ষর}
  \enspace$\bullet$\enspace
  \sfx{লাল = {\arabicfont \textarabic{علامة}} ও নতুন/বদলানো হরকত}
\end{center}
